\documentclass[12pt, a4paper]{ctexart}

% 常用宏包
\usepackage{geometry}
\geometry{left=2.5cm, right=2.5cm, top=2.5cm, bottom=2.5cm}
\usepackage{graphicx}
\usepackage{listings}
\usepackage{xcolor}
\usepackage{booktabs}
\usepackage{longtable}
\usepackage{hyperref}
\usepackage{float}
\usepackage{caption}
\usepackage{amsmath}

\hypersetup{
    colorlinks=true,
    linkcolor=blue,
    urlcolor=blue,
    citecolor=blue
}

% C 语言代码高亮样式配置
\definecolor{codegreen}{rgb}{0,0.6,0}
\definecolor{codegray}{rgb}{0.5,0.5,0.5}
\definecolor{codepurple}{rgb}{0.58,0,0.82}
\definecolor{backcolour}{rgb}{0.97,0.97,0.97}

\lstset{
    language=C,
    backgroundcolor=\color{backcolour},
    commentstyle=\color{codegreen},
    keywordstyle=\color{blue}\bfseries,
    numberstyle=\tiny\color{codegray},
    stringstyle=\color{codepurple},
    basicstyle=\ttfamily\small,
    breakatwhitespace=false,
    breaklines=true,
    captionpos=b,
    keepspaces=true,
    numbers=left,
    numbersep=5pt,
    showspaces=false,
    showstringspaces=false,
    showtabs=false,
    tabsize=2,
    frame=single,
    rulecolor=\color{gray!40}
}

\title{\textbf{MIT 6.1810: Operating System Engineering \\ Lab: Page Tables 实验报告}}
\author{姓名：XXX \quad 学号：XXXXXXXX}
\date{\today}

\begin{document}

\maketitle
\tableofcontents
\newpage

% ============================================================
\section{实验概述}
本实验围绕 RISC-V Sv39 分页机制展开，深入探索操作系统如何通过三级页表（Page Table）管理进程虚拟地址空间与物理内存映射。实验共包含以下四个任务：
\begin{enumerate}
    \item \textbf{Inspect a user-process page table (easy)}：运行 \texttt{pgtbltest}，针对 \texttt{print\_pgtbl} 输出的每一条页表项，解释它逻辑上对应进程的哪部分内存及权限位的含义。
    \item \textbf{Speed up system calls (easy)}：在进程创建时于 \texttt{USYSCALL} 处映射一页只读共享页，存入 \texttt{struct usyscall} 并初始化 \texttt{pid}，使 \texttt{ugetpid} 测试通过，并在进程退出时正确释放该页；同时回答还有哪些系统调用可以用共享页加速。
    \item \textbf{Print a page table (easy)}：在 \texttt{kernel/vm.c} 中实现 \texttt{vmprint()}，按官方说明的格式递归打印页表；并解释输出中各叶子页的内容、权限位，以及与 \texttt{print\_pgtbl} 输出的对应关系。
    \item \textbf{Use superpages (moderate / hard)}：修改内核使 \texttt{sbrk()} 申请含 2MB 对齐且不小于 2MB 的区域时使用单个超级页（Superpage）。实现 \texttt{superalloc()}/\texttt{superfree()}，修改 \texttt{uvmcopy()}、\texttt{uvmunmap()} 支持 \texttt{fork} 与释放（含超级页降级 Demote），使 \texttt{superpg\_fork}、\texttt{superpg\_free} 测试通过。
\end{enumerate}

% ============================================================
\section{任务一：Inspect a user-process page table}

\subsection{\texttt{print\_pgtbl} 原始输出与权限位原理}
在 RISC-V Sv39 分页方案中，虚拟地址低 39 位有效（xv6 为避免符号扩展将最高有效虚拟地址限制为 \texttt{MAXVA = 1L << 38 = 0x4000000000}），每页大小为 \texttt{PGSIZE = 4096} 字节（\texttt{0x1000}）。页表项（PTE）的高位（Bits 10..53）存放 44 位物理页号（PPN），低 10 位（Bits 0..9）为权限与状态标志位：
\begin{itemize}
    \item \texttt{PTE\_V (1 << 0 = 0x01)}：Valid，页表项有效位；
    \item \texttt{PTE\_R (1 << 1 = 0x02)}：Readable，可读权限位；
    \item \texttt{PTE\_W (1 << 2 = 0x04)}：Writable，可写权限位；
    \item \texttt{PTE\_X (1 << 3 = 0x08)}：Executable，可执行权限位；
    \item \texttt{PTE\_U (1 << 4 = 0x10)}：User，允许用户态（U-mode）访问（若未置位则仅 S-mode 内核态可访问）；
    \item \texttt{PTE\_A (1 << 6 = 0x40)}：Accessed，硬件自动置位，表示该页自上次清零后被读取或执行过；
    \item \texttt{PTE\_D (1 << 7 = 0x80)}：Dirty，硬件自动置位，表示该页自上次清零后被写入过。
\end{itemize}

运行 \texttt{pgtbltest} 时，\texttt{print\_pgtbl()} 输出了进程虚拟地址空间的前 10 页和最后 10 页（共 20 条 PTE）：
\begin{lstlisting}[language={}, caption={\texttt{print\_pgtbl} 终端输出}]
print_pgtbl starting
va 0x0 pte 0x21FC885B pa 0x87F22000 perm 0x5B
va 0x1000 pte 0x21FC7C5B pa 0x87F1F000 perm 0x5B
va 0x2000 pte 0x21FC7817 pa 0x87F1E000 perm 0x17
va 0x3000 pte 0x21FC7407 pa 0x87F1D000 perm 0x7
va 0x4000 pte 0x21FC70D7 pa 0x87F1C000 perm 0xD7
va 0x5000 pte 0x0 pa 0x0 perm 0x0
va 0x6000 pte 0x0 pa 0x0 perm 0x0
va 0x7000 pte 0x0 pa 0x0 perm 0x0
va 0x8000 pte 0x0 pa 0x0 perm 0x0
va 0x9000 pte 0x0 pa 0x0 perm 0x0
va 0x3FFFFF6000 pte 0x0 pa 0x0 perm 0x0
va 0x3FFFFF7000 pte 0x0 pa 0x0 perm 0x0
va 0x3FFFFF8000 pte 0x0 pa 0x0 perm 0x0
va 0x3FFFFF9000 pte 0x0 pa 0x0 perm 0x0
va 0x3FFFFFA000 pte 0x0 pa 0x0 perm 0x0
va 0x3FFFFFB000 pte 0x0 pa 0x0 perm 0x0
va 0x3FFFFFC000 pte 0x0 pa 0x0 perm 0x0
va 0x3FFFFFD000 pte 0x21FD4853 pa 0x87F52000 perm 0x53  (未实现任务二前为 0x0)
va 0x3FFFFFE000 pte 0x21FD08C7 pa 0x87F42000 perm 0xC7
va 0x3FFFFFF000 pte 0x2000184B pa 0x80006000 perm 0x4B
print_pgtbl: OK
\end{lstlisting}

\subsection{每一条页表项（PTE）的逐项详细解释}
针对上述 20 条输出，每一项在进程地址空间中对应的逻辑内容与权限位含义如表 \ref{tab:pgtbl-inspect} 所示：

\begin{longtable}{@{}llllp{6.0cm}@{}}
\caption{\texttt{print\_pgtbl} 全部 20 条页表项的逐项解析}
\label{tab:pgtbl-inspect} \\
\toprule
\textbf{虚拟地址 (\texttt{va})} & \textbf{物理地址 (\texttt{pa})} & \textbf{权限 (\texttt{perm})} & \textbf{权限位分解} & \textbf{逻辑上对应进程的内存区域与含义} \\
\midrule
\endfirsthead
\toprule
\textbf{虚拟地址 (\texttt{va})} & \textbf{物理地址 (\texttt{pa})} & \textbf{权限 (\texttt{perm})} & \textbf{权限位分解} & \textbf{逻辑上对应进程的内存区域与含义} \\
\midrule
\endhead
\texttt{0x0} & \texttt{0x87F22000} & \texttt{0x5B} & \texttt{A|U|X|R|V} & \textbf{用户代码段（Text 第 1 页）}：存放 \texttt{pgtbltest} 的机器指令与只读数据；用户态可读、可执行、不可写，已被 CPU 取指执行故 \texttt{A=1}。 \\
\texttt{0x1000} & \texttt{0x87F1F000} & \texttt{0x5B} & \texttt{A|U|X|R|V} & \textbf{用户代码段（Text 第 2 页）}：因代码超过 4KB 占用第 2 页。注意其物理地址与第 1 页不连续，因为 \texttt{kalloc()} 从空闲链表按 LIFO 顺序逐页分配。 \\
\texttt{0x2000} & \texttt{0x87F1E000} & \texttt{0x17} & \texttt{U|W|R|V} & \textbf{用户数据段（Data / BSS 页）}：存放已初始化和未初始化的全局/静态变量；用户态可读、可写、不可执行（无 \texttt{PTE\_X}）。 \\
\texttt{0x3000} & \texttt{0x87F1D000} & \texttt{0x07} & \texttt{W|R|V} & \textbf{用户栈保护页（Stack Guard Page）}：位于用户栈正下方，\texttt{kexec()} 通过 \texttt{uvmclear()} 清除了 \texttt{PTE\_U} 位。当用户栈向下溢出触碰该页时会因无用户态权限触发 Page Fault。 \\
\texttt{0x4000} & \texttt{0x87F1C000} & \texttt{0xD7} & \texttt{D|A|U|W|R|V} & \textbf{用户栈（User Stack 页）}：存放函数调用栈帧、局部变量及命令行参数 \texttt{argv}；用户态可读可写，因正在频繁读写栈内存，硬件已置位 \texttt{A=1, D=1}。 \\
\texttt{0x5000} & \texttt{0x0} & \texttt{0x0} & 无 (\texttt{PTE\_V=0}) & \textbf{未映射的堆空间}：初始进程大小 \texttt{p->sz = 0x5000}，尚未调用 \texttt{sbrk()} 扩展堆。 \\
\texttt{0x6000} & \texttt{0x0} & \texttt{0x0} & 无 (\texttt{PTE\_V=0}) & \textbf{未映射的堆空间}（无效 PTE）。 \\
\texttt{0x7000} & \texttt{0x0} & \texttt{0x0} & 无 (\texttt{PTE\_V=0}) & \textbf{未映射的堆空间}（无效 PTE）。 \\
\texttt{0x8000} & \texttt{0x0} & \texttt{0x0} & 无 (\texttt{PTE\_V=0}) & \textbf{未映射的堆空间}（无效 PTE）。 \\
\texttt{0x9000} & \texttt{0x0} & \texttt{0x0} & 无 (\texttt{PTE\_V=0}) & \textbf{未映射的堆空间}（无效 PTE）。 \\
\midrule
\texttt{0x3FFFFF6000} & \texttt{0x0} & \texttt{0x0} & 无 (\texttt{PTE\_V=0}) & \textbf{高地址未映射区域}（无效 PTE）。 \\
\texttt{0x3FFFFF7000} & \texttt{0x0} & \texttt{0x0} & 无 (\texttt{PTE\_V=0}) & \textbf{高地址未映射区域}（无效 PTE）。 \\
\texttt{0x3FFFFF8000} & \texttt{0x0} & \texttt{0x0} & 无 (\texttt{PTE\_V=0}) & \textbf{高地址未映射区域}（无效 PTE）。 \\
\texttt{0x3FFFFF9000} & \texttt{0x0} & \texttt{0x0} & 无 (\texttt{PTE\_V=0}) & \textbf{高地址未映射区域}（无效 PTE）。 \\
\texttt{0x3FFFFFA000} & \texttt{0x0} & \texttt{0x0} & 无 (\texttt{PTE\_V=0}) & \textbf{高地址未映射区域}（无效 PTE）。 \\
\texttt{0x3FFFFFB000} & \texttt{0x0} & \texttt{0x0} & 无 (\texttt{PTE\_V=0}) & \textbf{高地址未映射区域}（无效 PTE）。 \\
\texttt{0x3FFFFFC000} & \texttt{0x0} & \texttt{0x0} & 无 (\texttt{PTE\_V=0}) & \textbf{高地址未映射区域}（无效 PTE）。 \\
\texttt{0x3FFFFFD000} & \texttt{0x87F52000} & \texttt{0x13}/\texttt{0x53} & \texttt{(A)|U|R|V} & \textbf{\texttt{USYSCALL} 页（\texttt{TRAPFRAME - PGSIZE}）}：未实现任务二前为 \texttt{0x0}；实现后映射 \texttt{p->usyscall}，存放 \texttt{struct usyscall (pid)}，用户态只读共享（\texttt{PTE\_U | PTE\_R | PTE\_V}）。 \\
\texttt{0x3FFFFFE000} & \texttt{0x87F42000} & \texttt{0xC7} & \texttt{D|A|W|R|V} & \textbf{\texttt{TRAPFRAME} 页（\texttt{TRAMPOLINE - PGSIZE}）}：映射 \texttt{p->trapframe}，用于在进出内核时保存和恢复用户态的 32 个寄存器；可读可写但\textbf{无 \texttt{PTE\_U}}，仅内核态可访问。 \\
\texttt{0x3FFFFFF000} & \texttt{0x80006000} & \texttt{0x4B} & \texttt{A|X|R|V} & \textbf{\texttt{TRAMPOLINE} 页（\texttt{MAXVA - PGSIZE}）}：虚拟地址空间最高一页，映射内核中的 \texttt{trampoline.S} 跳板代码（\texttt{uservec}/\texttt{userret}）；可读可执行但\textbf{无 \texttt{PTE\_U}}，仅内核态可访问。 \\
\bottomrule
\end{longtable}

\begin{figure}[H]
    \centering
    % TODO: 替换为实际截图路径
    % \includegraphics[width=0.85\textwidth]{figs/print_pgtbl.png}
    \fbox{\parbox[c][4cm][c]{0.8\textwidth}{\centering \texttt{[图片占位符：\textbackslash includegraphics\{figs/print\_pgtbl.png\}]}}}
    \caption{\texttt{print\_pgtbl} 运行输出截图}
    \label{fig:print_pgtbl}
\end{figure}

% ============================================================
\section{任务二：Speed up system calls}

\subsection{设计思路}
为了加速 \texttt{getpid()} 系统调用，内核在每个进程创建时分配一页物理内存（\texttt{struct usyscall}），将进程的 \texttt{pid} 写入其中，并在进程页表中将其映射到虚拟地址 \texttt{USYSCALL}（\texttt{TRAPFRAME - PGSIZE = 0x3fffffd000}），权限设为 \texttt{PTE\_R | PTE\_U}（用户态只读）。这样用户态调用 \texttt{ugetpid()} 时可直接从虚拟地址 \texttt{USYSCALL} 读取 \texttt{pid}，省去了 \texttt{ecall} 陷入内核、保存上下文与页表切换的开销。

\subsection{核心代码实现}

\subsubsection{在进程控制块中添加 \texttt{usyscall} 指针（\texttt{kernel/proc.h}）}
\begin{lstlisting}[caption={\texttt{kernel/proc.h} 中 \texttt{struct proc} 新增字段}]
struct proc {
  ...
  uint64 sz;                   // Size of process memory (bytes)
  pagetable_t pagetable;       // User page table
  struct trapframe *trapframe; // data page for trampoline.S
  struct usyscall *usyscall;   // shared page for fast ugetpid()
  struct context context;      // swtch() here to run process
  ...
};
\end{lstlisting}

\subsubsection{分配、映射与释放 \texttt{usyscall} 页（\texttt{kernel/proc.c}）}
在 \texttt{allocproc()} 中分配物理页并初始化 \texttt{pid}；在 \texttt{proc\_pagetable()} 中建立 \texttt{USYSCALL} 映射（若映射失败需先解除已映射的 \texttt{TRAMPOLINE} 和 \texttt{TRAPFRAME}）；在 \texttt{proc\_freepagetable()} 和 \texttt{freeproc()} 中解除映射并回收物理页：

\begin{lstlisting}[caption={\texttt{kernel/proc.c} 中 \texttt{usyscall} 生命周期管理}]
static struct proc*
allocproc(void)
{
  ...
found:
  p->pid = allocpid();
  p->state = USED;

  // Allocate a trapframe page.
  if((p->trapframe = (struct trapframe *)kalloc()) == 0){
    freeproc(p);
    release(&p->lock);
    return 0;
  }

  // Allocate a usyscall page and initialize pid.
  if((p->usyscall = (struct usyscall *)kalloc()) == 0){
    freeproc(p);
    release(&p->lock);
    return 0;
  }
  p->usyscall->pid = p->pid;
  ...
}

static void
freeproc(struct proc *p)
{
  if(p->trapframe)
    kfree((void*)p->trapframe);
  p->trapframe = 0;
  if(p->usyscall)
    kfree((void*)p->usyscall);
  p->usyscall = 0;
  if(p->pagetable)
    proc_freepagetable(p->pagetable, p->sz);
  ...
}

pagetable_t
proc_pagetable(struct proc *p)
{
  ...
  if(mappages(pagetable, USYSCALL, PGSIZE,
              (uint64)(p->usyscall), PTE_R | PTE_U) < 0){
    uvmunmap(pagetable, TRAMPOLINE, 1, 0);
    uvmunmap(pagetable, TRAPFRAME, 1, 0);
    uvmfree(pagetable, 0);
    return 0;
  }

  return pagetable;
}

void
proc_freepagetable(pagetable_t pagetable, uint64 sz)
{
  uvmunmap(pagetable, TRAMPOLINE, 1, 0);
  uvmunmap(pagetable, TRAPFRAME, 1, 0);
  uvmunmap(pagetable, USYSCALL, 1, 0);
  uvmfree(pagetable, sz);
}
\end{lstlisting}

\subsection{思考题回答：还有哪些系统调用可以用共享页加速？}
凡是\textbf{仅从内核读取状态、不产生副作用且无需复杂权限校验的只读查询类系统调用}，都可以通过只读共享页（类似 Linux 内核的 \texttt{vvar}/\texttt{vDSO} 机制）进行加速：

\begin{enumerate}
    \item \textbf{\texttt{uptime()} / 时间类系统调用（如 Linux 中的 \texttt{gettimeofday()}、\texttt{clock\_gettime()}）}：
    \begin{itemize}
        \item \textbf{加速方式}：内核可以分配一页\textbf{全局共享物理页}，以只读权限（\texttt{PTE\_R | PTE\_U}）映射到所有用户进程的固定虚拟地址。每次时钟中断（\texttt{clockintr}）更新全局 \texttt{ticks} 时，内核同步更新该共享页中的计数值。用户态调用 \texttt{uptime()} 时直接读取该共享页即可，完全无需陷入内核。
    \end{itemize}
    \item \textbf{\texttt{getppid()}（获取父进程 ID）}：
    \begin{itemize}
        \item \textbf{加速方式}：在每个进程私有的 \texttt{struct usyscall} 页中增加 \texttt{int ppid;} 字段。在 \texttt{fork()} 创建子进程时写入父进程 PID，并在父进程退出导致子进程被托孤给 \texttt{init} 进程（\texttt{reparent()}）时由内核更新该字段，用户态即可零开销读取。
    \end{itemize}
    \item \textbf{进程身份与属性查询系统调用（如 \texttt{getuid()}、\texttt{geteuid()}、\texttt{getgid()}、\texttt{getpgrp()}）}：
    \begin{itemize}
        \item \textbf{加速方式}：将这些极少变动的进程属性一并存入进程的 \texttt{usyscall} 只读共享页，仅当进程调用 \texttt{setuid()} 等修改自身属性时由内核更新该页，所有查询调用均可在用户态直接完成。
    \end{itemize}
    \item \textbf{系统全局信息查询（如获取 CPU 核数、内存页大小、随机数种子或系统负载 \texttt{sysinfo}）}：
    \begin{itemize}
        \item \textbf{加速方式}：通过全局只读共享页向所有进程暴露静态硬件配置或由内核周期性更新的统计信息。
    \end{itemize}
\end{enumerate}

\subsection{测试结果}
\begin{figure}[H]
    \centering
    % TODO: 替换为实际截图路径
    % \includegraphics[width=0.85\textwidth]{figs/ugetpid_test.png}
    \fbox{\parbox[c][3.5cm][c]{0.8\textwidth}{\centering \texttt{[图片占位符：\textbackslash includegraphics\{figs/ugetpid\_test.png\}]}}}
    \caption{\texttt{ugetpid\_test} 运行结果截图}
    \label{fig:ugetpid_test}
\end{figure}

% ============================================================
\section{任务三：Print a page table}

\subsection{设计思路与核心代码实现}
\texttt{vmprint(pagetable\_t pagetable)} 用于按树形结构递归打印三级页表中的所有有效项：
\begin{itemize}
    \item 首先打印根页表（Level 2）的物理地址：\texttt{page table \%p}。
    \item 编写递归辅助函数 \texttt{vmprint\_walk(pagetable\_t pagetable, int level, uint64 va\_base)}，从顶层 \texttt{level = 2} 开始遍历当前页表的 512 个槽位。
    \item 若 \texttt{pte \& PTE\_V} 为真，打印 $3 - \text{level}$ 个 \texttt{" .."} 缩进，并通过 \texttt{va = va\_base | ((uint64)i << PXSHIFT(level))} 计算该 PTE 管辖的起始虚拟地址。
    \item 若该 PTE 的 \texttt{R/W/X} 位全为 0（即非叶子节点，指向下一级页表），则递归调用 \texttt{vmprint\_walk((pagetable\_t)PTE2PA(pte), level - 1, va)}。
\end{itemize}

\begin{lstlisting}[caption={\texttt{kernel/vm.c} 中 \texttt{vmprint} 与 \texttt{vmprint\_walk} 的实现}]
void
vmprint_walk(pagetable_t pagetable, int level, uint64 va_base)
{
  for(int i = 0; i < 512; i++){
    pte_t pte = pagetable[i];
    if((pte & PTE_V) != 0){
      for(int j = 0; j < 3 - level; j++)
        printf(" ..");

      uint64 va = ((uint64)i << PXSHIFT(level)) | va_base;
      uint64 child = PTE2PA(pte);
      printf("%p: pte %p pa %p\n", (void*)va, (void*)pte, (void*)child);
      if((pte & (PTE_R | PTE_W | PTE_X)) == 0){
        vmprint_walk((pagetable_t)PTE2PA(pte), level - 1, va);
      }
    }
  }
}

void
vmprint(pagetable_t pagetable)
{
  printf("page table %p\n", pagetable);
  vmprint_walk(pagetable, 2, 0);
}
\end{lstlisting}

\subsection{\texttt{print\_kpgtbl} 输出解析及与 \texttt{print\_pgtbl} 的对应关系}
运行 \texttt{print\_kpgtbl} 的实际输出如下：
\begin{lstlisting}[language={}, caption={\texttt{print\_kpgtbl} 终端输出}]
print_kpgtbl starting
page table 0x0000000087f22000
 ..0x0000000000000000: pte 0x0000000021fc7801 pa 0x0000000087f1e000
 .. ..0x0000000000000000: pte 0x0000000021fc7401 pa 0x0000000087f1d000
 .. .. ..0x0000000000000000: pte 0x0000000021fc7c5b pa 0x0000000087f1f000
 .. .. ..0x0000000000001000: pte 0x0000000021fc705b pa 0x0000000087f1c000
 .. .. ..0x0000000000002000: pte 0x0000000021fc6cd7 pa 0x0000000087f1b000
 .. .. ..0x0000000000003000: pte 0x0000000021fc6807 pa 0x0000000087f1a000
 .. .. ..0x0000000000004000: pte 0x0000000021fc64d7 pa 0x0000000087f19000
 ..0x0000003fc0000000: pte 0x0000000021fc8401 pa 0x0000000087f21000
 .. ..0x0000003fffe00000: pte 0x0000000021fc8001 pa 0x0000000087f20000
 .. .. ..0x0000003fffffd000: pte 0x0000000021fd4853 pa 0x0000000087f52000
 .. .. ..0x0000003fffffe000: pte 0x0000000021fd00c7 pa 0x0000000087f40000
 .. .. ..0x0000003ffffff000: pte 0x000000002000184b pa 0x0000000080006000
print_kpgtbl: OK
\end{lstlisting}

\subsubsection{1. 树形层级与中间页表项（非叶子节点）解析}
\begin{itemize}
    \item \textbf{第一级缩进（\texttt{" .."}，Level 2 顶层页表，每项跨度 $1\text{ GB} = \text{\texttt{0x40000000}}$）}：
    共有 2 项有效映射，权限位均为 \texttt{0x001}（仅 \texttt{PTE\_V=1}，指向下一级页表）：
    \begin{enumerate}
        \item \texttt{0x0000000000000000}（索引 \texttt{0}）：管辖低地址区 \texttt{[0x0, 0x40000000)}，指向物理地址为 \texttt{0x87f1e000} 的 Level-1 页表。
        \item \texttt{0x0000003fc0000000}（索引 \texttt{255 = 0xff}）：因 \texttt{MAXVA = 1L << 38}，最高位第 38 位恒为 0，故顶层最大索引为 $2^8 - 1 = 255$（\texttt{255 << 30 = 0x3fc0000000}），管辖高地址区 \texttt{[0x3fc0000000, 0x4000000000)}，指向物理地址为 \texttt{0x87f21000} 的 Level-1 页表。
    \end{enumerate}
    \item \textbf{第二级缩进（\texttt{" .. .."}，Level 1 中间页表，每项跨度 $2\text{ MB} = \text{\texttt{0x200000}}$）}：
    权限位同样为 \texttt{0x001}（指向 Level-0 页表）：
    \begin{enumerate}
        \item \texttt{0x0000000000000000}（L2 索引 0 下的 L1 索引 \texttt{0}）：管辖 \texttt{[0x0, 0x200000)}，指向 \texttt{0x87f1d000} 的 Level-0 页表。
        \item \texttt{0x0000003fffe00000}（L2 索引 255 下的 L1 索引 \texttt{511}）：管辖最高 2MB 空间 \texttt{[0x3fffe00000, 0x4000000000)}，指向 \texttt{0x87f20000} 的 Level-0 页表。
    \end{enumerate}
\end{itemize}

\subsubsection{2. 各叶子页（\texttt{" .. .. .."}，Level 0）的内容、权限位及与 \texttt{print\_pgtbl} 的对应关系}
\texttt{vmprint} 只打印有效（\texttt{PTE\_V = 1}）的页表项，而 \texttt{print\_pgtbl} 則是调用 \texttt{pgpte(va)} 枚举前 10 页和后 10 页（包含未映射的 \texttt{0x0} 项）。两者之间的对应关系是：\textbf{\texttt{vmprint} 打印出的 8 个三级缩进（\texttt{" .. .. .."}）叶子节点，与 \texttt{print\_pgtbl} 输出中所有非零的有效页表项呈严格的一一对应关系}（仅在 \texttt{print\_kpgtbl} 运行前由于 \texttt{ugetpid\_test} 访问了全局变量和 \texttt{USYSCALL} 页，使得 \texttt{0x2000} 数据段页的 \texttt{D|A} 位以及 \texttt{0x3fffffd000} \texttt{USYSCALL} 页的 \texttt{A} 位被硬件置位）：

\begin{table}[H]
\centering
\small
\caption{\texttt{vmprint} 叶子节点与 \texttt{print\_pgtbl} 输出的对应关系表}
\label{tab:vmprint-leaf}
\begin{tabular}{@{}llllp{5.5cm}@{}}
\toprule
\textbf{\texttt{vmprint} 叶子虚拟地址} & \textbf{L2/L1/L0 索引} & \textbf{权限位 (\texttt{perm})} & \textbf{对应 \texttt{print\_pgtbl} 项} & \textbf{叶子页逻辑内容与权限说明} \\
\midrule
\texttt{0x0000000000000000} & \texttt{0 / 0 / 0} & \texttt{0x5b (A|U|X|R|V)} & \texttt{va 0x0} & 用户代码段第 1 页（只读可执行） \\
\texttt{0x0000000000001000} & \texttt{0 / 0 / 1} & \texttt{0x5b (A|U|X|R|V)} & \texttt{va 0x1000} & 用户代码段第 2 页（只读可执行） \\
\texttt{0x0000000000002000} & \texttt{0 / 0 / 2} & \texttt{0xd7 (D|A|U|W|R|V)} & \texttt{va 0x2000} & 用户数据段（\texttt{ugetpid\_test} 修改了全局变量 \texttt{testname}，故相比 \texttt{print\_pgtbl} 时的 \texttt{0x17} 多了 \texttt{D|A} 位） \\
\texttt{0x0000000000003000} & \texttt{0 / 0 / 3} & \texttt{0x07 (W|R|V)} & \texttt{va 0x3000} & 栈保护页（Guard Page，无 \texttt{PTE\_U}） \\
\texttt{0x0000000000004000} & \texttt{0 / 0 / 4} & \texttt{0xd7 (D|A|U|W|R|V)} & \texttt{va 0x4000} & 用户栈页（可读可写，已访问且脏） \\
\midrule
\texttt{0x0000003fffffd000} & \texttt{255 / 511 / 509} & \texttt{0x53 (A|U|R|V)} & \texttt{va 0x3FFFFFD000} & \texttt{USYSCALL} 共享页（用户态只读，因 \texttt{ugetpid()} 已读取过故 \texttt{A=1}） \\
\texttt{0x0000003fffffe000} & \texttt{255 / 511 / 510} & \texttt{0xc7 (D|A|W|R|V)} & \texttt{va 0x3FFFFFE000} & \texttt{TRAPFRAME} 页（内核态读写，无 \texttt{PTE\_U}） \\
\texttt{0x0000003ffffff000} & \texttt{255 / 511 / 511} & \texttt{0x4b (A|X|R|V)} & \texttt{va 0x3FFFFFF000} & \texttt{TRAMPOLINE} 跳板代码页（内核态只读可执行，无 \texttt{PTE\_U}） \\
\bottomrule
\end{tabular}
\end{table}

而 \texttt{print\_pgtbl} 中其余值为 \texttt{0x0} 的虚拟地址（\texttt{0x5000}$\sim$\texttt{0x9000} 以及 \texttt{0x3FFFFF6000}$\sim$\texttt{0x3FFFFFC000}）因为尚未映射（\texttt{PTE\_V = 0}），在 \texttt{vmprint} 中被自动跳过不打印。

\begin{figure}[H]
    \centering
    % TODO: 替换为实际截图路径
    % \includegraphics[width=0.85\textwidth]{figs/print_kpgtbl.png}
    \fbox{\parbox[c][4cm][c]{0.8\textwidth}{\centering \texttt{[图片占位符：\textbackslash includegraphics\{figs/print\_kpgtbl.png\}]}}}
    \caption{\texttt{print\_kpgtbl} 页表打印截图}
    \label{fig:print_kpgtbl}
\end{figure}

% ============================================================
\section{任务四：Use superpages}

\subsection{设计思路}
RISC-V Sv39 支持在 Level-1 页表项中直接设置 \texttt{PTE\_V | PTE\_R/W/X}，从而将单个 L1 PTE 映射为 2MB（\texttt{SUPERPGSIZE = 2 * (1 << 20)}）的物理超级页（Megapage / Superpage）。为实现完整的超级页支持，设计分为五个关键部分：
\begin{enumerate}
    \item \textbf{双空闲链表物理内存分配器（\texttt{kernel/kalloc.c}）}：新增独立的 \texttt{supermem} 空闲链表，在 \texttt{freerange()} 初始化时预留靠近 \texttt{PHYSTOP} 的 16 个 2MB 对齐物理块放入 \texttt{supermem.freelist}，并实现 \texttt{superalloc()} 和 \texttt{superfree()}。
    \item \textbf{大页标志位与 L1 映射（\texttt{kernel/riscv.h}, \texttt{kernel/vm.c}）}：利用 RISC-V PTE 的软件保留位（RSW）定义 \texttt{\#define PTE\_S (1L << 8)} 标识超级页；实现只遍历到 Level-1 的 \texttt{superwalk()} 和 \texttt{mapsuperpages()}。
    \item \textbf{按需混合分配（\texttt{uvmalloc}）}：在 \texttt{uvmalloc()} 循环中，当虚拟地址 \texttt{a} 满足 2MB 对齐（\texttt{a \% SUPERPGSIZE == 0}）且剩余待分配大小 $\ge$ 2MB（\texttt{a + SUPERPGSIZE <= newsz}）时，优先调用 \texttt{superalloc()} 与 \texttt{mapsuperpages()} 建立超级页映射；其余零头区域或超级页耗尽时使用 4KB 普通页。
    \item \textbf{支持超级页的进程复制（\texttt{uvmcopy}）}：\texttt{fork()} 时根据父进程 PTE 是否带有 \texttt{PTE\_S}，动态调整步长 \texttt{szinc} 为 \texttt{SUPERPGSIZE} 或 \texttt{PGSIZE}，调用 \texttt{superalloc()} 与 \texttt{mapsuperpages()} 复制 2MB 超级页。
    \item \textbf{超级页释放与部分释放降级（\texttt{uvmunmap} \& Demote）}：
    \begin{itemize}
        \item 若待解除映射范围完整覆盖整个 2MB 超级页，则直接调用 \texttt{superfree()} 并清空 L1 PTE；
        \item 若只释放超级页的一部分（如 \texttt{sbrk(-PGSIZE)}），则触发\textbf{超级页降级（Demote）}：先清空原 L1 PTE，分配 512 个 4KB 普通页，将原 2MB 超级页内容逐页拷贝并映射到 Level-0 页表，释放原 2MB 物理超级页，最后再按普通页释放目标 4KB 页。
    \end{itemize}
\end{enumerate}

\subsection{核心代码实现}

\subsubsection{2MB 物理超级页分配器（\texttt{kernel/kalloc.c}）}
\begin{lstlisting}[caption={\texttt{kernel/kalloc.c} 中预留 2MB 超级页与 \texttt{superalloc}/\texttt{superfree} 实现}]
struct {
  struct spinlock lock;
  struct run *freelist;
} kmem, supermem;

void
kinit()
{
  initlock(&kmem.lock, "kmem");
  initlock(&supermem.lock, "supermem");
  freerange(end, (void*)PHYSTOP);
}

void
freerange(void *pa_start, void *pa_end)
{
  char *p;
  char *super_start = (char*)pa_end - 16 * SUPERPGSIZE;

  p = (char*)PGROUNDUP((uint64)pa_start);
  for(; p + PGSIZE <= (char*)super_start; p += PGSIZE)
    kfree(p);

  for(p = (char*)SUPERPGROUNDUP((uint64)super_start); p + SUPERPGSIZE <= (char*)pa_end; p += SUPERPGSIZE)
    superfree(p);
}

void *
superalloc(void)
{
  struct run *r;

  acquire(&supermem.lock);
  r = supermem.freelist;
  if(r)
    supermem.freelist = r->next;
  release(&supermem.lock);

  if(r)
    memset((char*)r, 5, SUPERPGSIZE); // fill with junk
  return (void*)r;
}

void
superfree(void *pa)
{
  struct run *r;

  if(((uint64)pa % SUPERPGSIZE) != 0 || (char*)pa < end || (uint64)pa >= PHYSTOP)
    panic("superfree");

  memset(pa, 1, SUPERPGSIZE);
  r = (struct run*)pa;

  acquire(&supermem.lock);
  r->next = supermem.freelist;
  supermem.freelist = r;
  release(&supermem.lock);
}
\end{lstlisting}

\subsubsection{L1 超级页查找与映射（\texttt{kernel/vm.c}）}
\begin{lstlisting}[caption={\texttt{superwalk}、\texttt{mapsuperpages} 与 \texttt{walkaddr} 实现}]
pte_t *
superwalk(pagetable_t pagetable, uint64 va, int alloc)
{
  if(va >= MAXVA)
    panic("superwalk");
  pte_t *pte = &pagetable[PX(2, va)];
  if(*pte & PTE_V) {
    pagetable = (pagetable_t)PTE2PA(*pte);
  } else {
    if(!alloc || (pagetable = (pde_t*)kalloc()) == 0)
      return 0;
    memset(pagetable, 0, PGSIZE);
    *pte = PA2PTE(pagetable) | PTE_V;
  }
  return &pagetable[PX(1, va)];
}

int
mapsuperpages(pagetable_t pagetable, uint64 va, uint64 size, uint64 pa, int perm)
{
  uint64 a, last;
  pte_t *pte;

  if((va % SUPERPGSIZE) != 0)
    panic("mapsuperpages: va not aligned");
  if((size % SUPERPGSIZE) != 0)
    panic("mapsuperpages: size not aligned");
  if(size == 0)
    panic("mapsuperpages: size");

  a = va;
  last = va + size - SUPERPGSIZE;
  for(;;){
    if((pte = superwalk(pagetable, a, 1)) == 0)
      return -1;
    if(*pte & PTE_V)
      panic("mapsuperpages: remap");
    *pte = PA2PTE(pa) | perm | PTE_V | PTE_S;
    if(a == last)
      break;
    a += SUPERPGSIZE;
    pa += SUPERPGSIZE;
  }
  return 0;
}
\end{lstlisting}

\subsubsection{超级页分配、复制与降级释放（\texttt{uvmalloc}, \texttt{uvmcopy}, \texttt{uvmunmap}）}
\begin{lstlisting}[caption={\texttt{uvmunmap}（含超级页降级 Demote）、\texttt{uvmalloc} 与 \texttt{uvmcopy} 核心实现}]
void
uvmunmap(pagetable_t pagetable, uint64 va, uint64 npages, int do_free)
{
  uint64 a;
  pte_t *pte;
  int sz = PGSIZE;

  if((va % PGSIZE) != 0)
    panic("uvmunmap: not aligned");

  for(a = va; a < va + npages*PGSIZE; a += sz){
    sz = PGSIZE;
    if((pte = walk(pagetable, a, 0)) == 0)
      continue;
    if((*pte & PTE_V) == 0)
      continue;

    if(*pte & PTE_S){
      // 情况 1：整块 2MB 大页都在本次解除映射的范围内
      if(a % SUPERPGSIZE == 0 && a + SUPERPGSIZE <= va + npages*PGSIZE){
        sz = SUPERPGSIZE;
        if(do_free){
          uint64 pa = PTE2PA(*pte);
          superfree((void*)pa);
        }
        *pte = 0;
        continue;
      } else {
        // 情况 2：只释放大页的一部分（Demote 降级为 512 个 4KB 普通页）
        uint64 super_pa = PTE2PA(*pte);
        int flags = PTE_FLAGS(*pte) & ~(PTE_S | PTE_V);
        uint64 s = SUPERPGROUNDDOWN(a);

        *pte = 0;
        for(uint64 pa_off = 0; pa_off < SUPERPGSIZE; pa_off += PGSIZE){
          char *mem = kalloc();
          if(mem == 0)
            panic("uvmunmap: demote out of memory");
          memmove(mem, (char*)(super_pa + pa_off), PGSIZE);
          if(mappages(pagetable, s + pa_off, PGSIZE, (uint64)mem, flags) != 0)
            panic("uvmunmap: demote mappages");
        }
        superfree((void*)super_pa);
        pte = walk(pagetable, a, 0);
      }
    }

    sz = PGSIZE;
    if(PTE_FLAGS(*pte) == PTE_V)
      panic("uvmunmap: not a leaf");
    if(do_free){
      uint64 pa = PTE2PA(*pte);
      kfree((void*)pa);
    }
    *pte = 0;
  }
}

uint64
uvmalloc(pagetable_t pagetable, uint64 oldsz, uint64 newsz, int xperm)
{
  char *mem;
  uint64 a;
  int sz;

  if(newsz < oldsz)
    return oldsz;

  oldsz = PGROUNDUP(oldsz);
  for(a = oldsz; a < newsz; a += sz){
    if(a % SUPERPGSIZE == 0 && a + SUPERPGSIZE <= newsz && (mem = superalloc()) != 0){
      sz = SUPERPGSIZE;
      #ifndef LAB_SYSCALL
      memset(mem, 0, sz);
      #endif
      if(mapsuperpages(pagetable, a, sz, (uint64)mem, PTE_R|PTE_U|xperm) != 0){
        superfree(mem);
        uvmdealloc(pagetable, a, oldsz);
        return 0;
      }
    } else {
      sz = PGSIZE;
      mem = kalloc();
      if(mem == 0){
        uvmdealloc(pagetable, a, oldsz);
        return 0;
      }
      #ifndef LAB_SYSCALL
      memset(mem, 0, sz);
      #endif
      if(mappages(pagetable, a, sz, (uint64)mem, PTE_R|PTE_U|xperm) != 0){
        kfree(mem);
        uvmdealloc(pagetable, a, oldsz);
        return 0;
      }
    }
  }
  return newsz;
}

int
uvmcopy(pagetable_t old, pagetable_t new, uint64 sz)
{
  pte_t *pte;
  uint64 pa, i;
  uint flags;
  char *mem;
  int szinc = PGSIZE;

  for(i = 0; i < sz; i += szinc){
    if((pte = walk(old, i, 0)) == 0)
      continue;
    if((*pte & PTE_V) == 0)
      continue;

    if(*pte & PTE_S){
      szinc = SUPERPGSIZE;
      pa = PTE2PA(*pte);
      flags = PTE_FLAGS(*pte);
      if((mem = superalloc()) == 0)
        goto err;
      memmove(mem, (char*)pa, SUPERPGSIZE);
      if(mapsuperpages(new, i, SUPERPGSIZE, (uint64)mem, flags) != 0){
        superfree(mem);
        goto err;
      }
    } else {
      szinc = PGSIZE;
      pa = PTE2PA(*pte);
      flags = PTE_FLAGS(*pte);
      if((mem = kalloc()) == 0)
        goto err;
      memmove(mem, (char*)pa, PGSIZE);
      if(mappages(new, i, PGSIZE, (uint64)mem, flags) != 0){
        kfree(mem);
        goto err;
      }
    }
  }
  return 0;

 err:
  uvmunmap(new, 0, i / PGSIZE, 1);
  return -1;
}
\end{lstlisting}

\subsection{测试结果}
运行 \texttt{pgtbltest}，\texttt{superpg\_fork} 与 \texttt{superpg\_free} 测试均顺利通过：
\begin{lstlisting}[language={}, caption={\texttt{superpg\_fork} 与 \texttt{superpg\_free} 运行输出}]
superpg_fork starting
usertrap(): unexpected scause 0xf pid=69
            sepc=0x2ec stval=0x5001
superpg_fork: OK
superpg_free starting
usertrap(): unexpected scause 0xd pid=70
            sepc=0x40a stval=0xfff001
superpg_free: OK
pgtbltest: all tests succeeded
\end{lstlisting}

\begin{figure}[H]
    \centering
    % TODO: 替换为实际截图路径
    % \includegraphics[width=0.85\textwidth]{figs/superpg_test.png}
    \fbox{\parbox[c][4cm][c]{0.8\textwidth}{\centering \texttt{[图片占位符：\textbackslash includegraphics\{figs/superpg\_test.png\}]}}}
    \caption{\texttt{superpg\_fork} 与 \texttt{superpg\_free} 测试通过截图}
    \label{fig:superpg_test}
\end{figure}

% ============================================================
\section{综合测试与 \texttt{make grade} 评测结果}

在 QEMU 中运行 \texttt{usertests -q} 验证所有回归测试，随后在宿主机终端执行 \texttt{make grade} 进行自动评测，所有测试项均顺利通过。

\begin{figure}[H]
    \centering
    % TODO: 替换为实际截图路径
    % \includegraphics[width=0.85\textwidth]{figs/usertests.png}
    \fbox{\parbox[c][4cm][c]{0.8\textwidth}{\centering \texttt{[图片占位符：\textbackslash includegraphics\{figs/usertests.png\}]}}}
    \caption{\texttt{usertests -q} 回归测试通过截图}
    \label{fig:usertests}
\end{figure}

\begin{figure}[H]
    \centering
    % TODO: 替换为实际截图路径
    % \includegraphics[width=0.85\textwidth]{figs/make_grade.png}
    \fbox{\parbox[c][4.5cm][c]{0.8\textwidth}{\centering \texttt{[图片占位符：\textbackslash includegraphics\{figs/make\_grade.png\}]}}}
    \caption{\texttt{make grade} 满分通过截图}
    \label{fig:make_grade}
\end{figure}

\end{document}
